\IfFileExists{stacks-project.cls}{%
\documentclass{stacks-project}
}{%
\documentclass{amsart}
}

% For dealing with references we use the comment environment
\usepackage{verbatim}
\newenvironment{reference}{\comment}{\endcomment}
%\newenvironment{reference}{}{}
\newenvironment{slogan}{\comment}{\endcomment}
\newenvironment{history}{\comment}{\endcomment}

% For commutative diagrams we use Xy-pic
\usepackage[all]{xy}

% We use 2cell for 2-commutative diagrams.
\xyoption{2cell}
\UseAllTwocells

% We use multicol for the list of chapters between chapters
\usepackage{multicol}

% This is generally recommended for better output
\usepackage{lmodern}
\usepackage[T1]{fontenc}

% For cross-file-references
\usepackage{xr-hyper}

% Package for hypertext links:
\usepackage{hyperref}

% For any local file, say "hello.tex" you want to link to please
% use \externaldocument[hello-]{hello}
\externaldocument[introduction-]{introduction}
\externaldocument[conventions-]{conventions}
\externaldocument[sets-]{sets}
\externaldocument[categories-]{categories}
\externaldocument[topology-]{topology}
\externaldocument[sheaves-]{sheaves}
\externaldocument[sites-]{sites}
\externaldocument[stacks-]{stacks}
\externaldocument[fields-]{fields}
\externaldocument[algebra-]{algebra}
\externaldocument[brauer-]{brauer}
\externaldocument[homology-]{homology}
\externaldocument[derived-]{derived}
\externaldocument[simplicial-]{simplicial}
\externaldocument[more-algebra-]{more-algebra}
\externaldocument[smoothing-]{smoothing}
\externaldocument[modules-]{modules}
\externaldocument[sites-modules-]{sites-modules}
\externaldocument[injectives-]{injectives}
\externaldocument[cohomology-]{cohomology}
\externaldocument[sites-cohomology-]{sites-cohomology}
\externaldocument[dga-]{dga}
\externaldocument[dpa-]{dpa}
\externaldocument[sdga-]{sdga}
\externaldocument[hypercovering-]{hypercovering}
\externaldocument[schemes-]{schemes}
\externaldocument[constructions-]{constructions}
\externaldocument[properties-]{properties}
\externaldocument[morphisms-]{morphisms}
\externaldocument[coherent-]{coherent}
\externaldocument[divisors-]{divisors}
\externaldocument[limits-]{limits}
\externaldocument[varieties-]{varieties}
\externaldocument[topologies-]{topologies}
\externaldocument[descent-]{descent}
\externaldocument[perfect-]{perfect}
\externaldocument[more-morphisms-]{more-morphisms}
\externaldocument[flat-]{flat}
\externaldocument[groupoids-]{groupoids}
\externaldocument[more-groupoids-]{more-groupoids}
\externaldocument[etale-]{etale}
\externaldocument[chow-]{chow}
\externaldocument[intersection-]{intersection}
\externaldocument[pic-]{pic}
\externaldocument[weil-]{weil}
\externaldocument[adequate-]{adequate}
\externaldocument[dualizing-]{dualizing}
\externaldocument[duality-]{duality}
\externaldocument[discriminant-]{discriminant}
\externaldocument[derham-]{derham}
\externaldocument[local-cohomology-]{local-cohomology}
\externaldocument[algebraization-]{algebraization}
\externaldocument[curves-]{curves}
\externaldocument[resolve-]{resolve}
\externaldocument[models-]{models}
\externaldocument[functors-]{functors}
\externaldocument[equiv-]{equiv}
\externaldocument[pione-]{pione}
\externaldocument[etale-cohomology-]{etale-cohomology}
\externaldocument[proetale-]{proetale}
\externaldocument[relative-cycles-]{relative-cycles}
\externaldocument[more-etale-]{more-etale}
\externaldocument[trace-]{trace}
\externaldocument[crystalline-]{crystalline}
\externaldocument[spaces-]{spaces}
\externaldocument[spaces-properties-]{spaces-properties}
\externaldocument[spaces-morphisms-]{spaces-morphisms}
\externaldocument[decent-spaces-]{decent-spaces}
\externaldocument[spaces-cohomology-]{spaces-cohomology}
\externaldocument[spaces-limits-]{spaces-limits}
\externaldocument[spaces-divisors-]{spaces-divisors}
\externaldocument[spaces-over-fields-]{spaces-over-fields}
\externaldocument[spaces-topologies-]{spaces-topologies}
\externaldocument[spaces-descent-]{spaces-descent}
\externaldocument[spaces-perfect-]{spaces-perfect}
\externaldocument[spaces-more-morphisms-]{spaces-more-morphisms}
\externaldocument[spaces-flat-]{spaces-flat}
\externaldocument[spaces-groupoids-]{spaces-groupoids}
\externaldocument[spaces-more-groupoids-]{spaces-more-groupoids}
\externaldocument[bootstrap-]{bootstrap}
\externaldocument[spaces-pushouts-]{spaces-pushouts}
\externaldocument[spaces-chow-]{spaces-chow}
\externaldocument[groupoids-quotients-]{groupoids-quotients}
\externaldocument[spaces-more-cohomology-]{spaces-more-cohomology}
\externaldocument[spaces-simplicial-]{spaces-simplicial}
\externaldocument[spaces-duality-]{spaces-duality}
\externaldocument[formal-spaces-]{formal-spaces}
\externaldocument[restricted-]{restricted}
\externaldocument[spaces-resolve-]{spaces-resolve}
\externaldocument[formal-defos-]{formal-defos}
\externaldocument[defos-]{defos}
\externaldocument[cotangent-]{cotangent}
\externaldocument[examples-defos-]{examples-defos}
\externaldocument[algebraic-]{algebraic}
\externaldocument[examples-stacks-]{examples-stacks}
\externaldocument[stacks-sheaves-]{stacks-sheaves}
\externaldocument[criteria-]{criteria}
\externaldocument[artin-]{artin}
\externaldocument[quot-]{quot}
\externaldocument[stacks-properties-]{stacks-properties}
\externaldocument[stacks-morphisms-]{stacks-morphisms}
\externaldocument[stacks-limits-]{stacks-limits}
\externaldocument[stacks-cohomology-]{stacks-cohomology}
\externaldocument[stacks-perfect-]{stacks-perfect}
\externaldocument[stacks-introduction-]{stacks-introduction}
\externaldocument[stacks-more-morphisms-]{stacks-more-morphisms}
\externaldocument[stacks-geometry-]{stacks-geometry}
\externaldocument[moduli-]{moduli}
\externaldocument[moduli-curves-]{moduli-curves}
\externaldocument[examples-]{examples}
\externaldocument[exercises-]{exercises}
\externaldocument[guide-]{guide}
\externaldocument[desirables-]{desirables}
\externaldocument[coding-]{coding}
\externaldocument[obsolete-]{obsolete}
\externaldocument[fdl-]{fdl}
\externaldocument[index-]{index}

% Theorem environments.
%
\theoremstyle{plain}
\newtheorem{theorem}[subsection]{Theorem}
\newtheorem{proposition}[subsection]{Proposition}
\newtheorem{lemma}[subsection]{Lemma}

\theoremstyle{definition}
\newtheorem{definition}[subsection]{Definition}
\newtheorem{example}[subsection]{Example}
\newtheorem{exercise}[subsection]{Exercise}
\newtheorem{situation}[subsection]{Situation}

\theoremstyle{remark}
\newtheorem{remark}[subsection]{Remark}
\newtheorem{remarks}[subsection]{Remarks}

\numberwithin{equation}{subsection}

% Macros
%
\def\lim{\mathop{\mathrm{lim}}\nolimits}
\def\colim{\mathop{\mathrm{colim}}\nolimits}
\def\Spec{\mathop{\mathrm{Spec}}}
\def\Hom{\mathop{\mathrm{Hom}}\nolimits}
\def\Ext{\mathop{\mathrm{Ext}}\nolimits}
\def\SheafHom{\mathop{\mathcal{H}\!\mathit{om}}\nolimits}
\def\SheafExt{\mathop{\mathcal{E}\!\mathit{xt}}\nolimits}
\def\Sch{\mathit{Sch}}
\def\Mor{\mathop{\mathrm{Mor}}\nolimits}
\def\Ob{\mathop{\mathrm{Ob}}\nolimits}
\def\Sh{\mathop{\mathit{Sh}}\nolimits}
\def\NL{\mathop{N\!L}\nolimits}
\def\CH{\mathop{\mathrm{CH}}\nolimits}
\def\proetale{{pro\text{-}\acute{e}tale}}
\def\etale{{\acute{e}tale}}
\def\QCoh{\mathit{QCoh}}
\def\Ker{\mathop{\mathrm{Ker}}}
\def\Im{\mathop{\mathrm{Im}}}
\def\Coker{\mathop{\mathrm{Coker}}}
\def\Coim{\mathop{\mathrm{Coim}}}

% Boxtimes
%
\DeclareMathSymbol{\boxtimes}{\mathbin}{AMSa}{"02}

%
% Macros for moduli stacks/spaces
%
\def\QCohstack{\mathcal{QC}\!\mathit{oh}}
\def\Cohstack{\mathcal{C}\!\mathit{oh}}
\def\Spacesstack{\mathcal{S}\!\mathit{paces}}
\def\Quotfunctor{\mathrm{Quot}}
\def\Hilbfunctor{\mathrm{Hilb}}
\def\Curvesstack{\mathcal{C}\!\mathit{urves}}
\def\Polarizedstack{\mathcal{P}\!\mathit{olarized}}
\def\Complexesstack{\mathcal{C}\!\mathit{omplexes}}
% \Pic is the operator that assigns to X its picard group, usage \Pic(X)
% \Picardstack_{X/B} denotes the Picard stack of X over B
% \Picardfunctor_{X/B} denotes the Picard functor of X over B
\def\Pic{\mathop{\mathrm{Pic}}\nolimits}
\def\Picardstack{\mathcal{P}\!\mathit{ic}}
\def\Picardfunctor{\mathrm{Pic}}
\def\Deformationcategory{\mathcal{D}\!\mathit{ef}}
\usepackage{cleveref}
\crefname{originalreference}{original reference}{original references}
\Crefname{originalreference}{Original reference}{Original references}
\makeatletter
\newlabel{categories-lemma-representable-diagonal-groupoids}{{02YA}{source}{}{sourcepointer-1}{}}
\newlabel{categories-lemma-representable-diagonal-groupoids@cref}{{[originalreference][0][]02YA}{[1][0][]source}{}{}{}}
\makeatother
\crefname{originallemma}{original lemma}{original lemmas}
\Crefname{originallemma}{Original lemma}{Original lemmas}
\makeatletter
\newlabel{lemma-base-change-by-space-representable-by-space}{{0300}{source}{}{sourcepointer-2}{}}
\newlabel{lemma-base-change-by-space-representable-by-space@cref}{{[originallemma][0][]0300}{[1][0][]source}{}{}{}}
\makeatother
\crefname{originallemma}{original lemma}{original lemmas}
\Crefname{originallemma}{Original lemma}{Original lemmas}
\makeatletter
\newlabel{lemma-representable-by-spaces-morphism-equivalent}{{0457}{source}{}{sourcepointer-3}{}}
\newlabel{lemma-representable-by-spaces-morphism-equivalent@cref}{{[originallemma][0][]0457}{[1][0][]source}{}{}{}}
\makeatother
\crefname{originalreference}{original reference}{original references}
\Crefname{originalreference}{Original reference}{Original references}
\makeatletter
\newlabel{stacks-lemma-isom-as-2-fibre-product}{{04SI}{source}{}{sourcepointer-4}{}}
\newlabel{stacks-lemma-isom-as-2-fibre-product@cref}{{[originalreference][0][]04SI}{[1][0][]source}{}{}{}}
\makeatother
\title{Eight equivalent characterizations of representability of the diagonal by algebraic spaces}
\author{The Stacks Project Authors}
\date{stacks-project:algebraic.tex:local-snapshot-20261012; GNU FDL 1.2 or later}
\begin{document}
\maketitle
\noindent\textit{Editorial note.} The three 'formal' pairs are restatements of the same quantified statement at different representation levels; each direction is carried as an explicit obligation with the identity spelled out. The equivalence-invariance input is not carded ('Omitted.') and is declared external.\par\medskip
\noindent\textit{Allowed standard prerequisites (editorial).}\begin{itemize}
\item Isom sheaf as 2-fibre product. Stacks, lemma isom-as-2-fibre-product (cross-chapter, named). Gives (1) iff (2).
\item 2-fibre product algebra of categories. Categories, lemmas 2-product-categories-over-C, representable-diagonal-groupoids (cross-chapter, named). The rewriting identities used in (1) iff (3).
\item Equivalence invariance of representability. lemma-representable-by-spaces-morphism-equivalent (same chapter, NOT carded, declared external). Gives (3) iff (5) and (4) iff (6).
\item Base change by a space-representable category. lemma-base-change-by-space-representable-by-space (carded L26677). Gives (3) implies (7).
\item Fibre products of algebraic spaces. Spaces, lemma fibre-product-spaces (cross-chapter, named). Makes the rewritten right-hand side representable.
\end{itemize}\medskip
\ifdefined\hypertarget\hypertarget{sourcepointer-1}{}\fi
\noindent\textit{Source reference.} Original reference 02YA: focused-seeds/geometry/S9-alg-representable-diagonal/ref-tags.txt, source lines 5--5. Navigation pointer only: the counted proof cites this statement by \textbackslash{}ref as a standing input; the deduction the counted proof has to make is carried out inside the counted span, and no part of that deduction is delegated to this pointer.\par\medskip
\ifdefined\hypertarget\hypertarget{sourcepointer-2}{}\fi
\noindent\textit{Source reference.} Original lemma 0300: sources/stacks-algebraic/pinned-original-source.tex, source lines 900--900. Navigation pointer only: the counted proof cites this statement by \textbackslash{}ref as a standing input; the deduction the counted proof has to make is carried out inside the counted span, and no part of that deduction is delegated to this pointer.\par\medskip
\ifdefined\hypertarget\hypertarget{sourcepointer-3}{}\fi
\noindent\textit{Source reference.} Original lemma 0457: sources/stacks-algebraic/pinned-original-source.tex, source lines 802--802. Navigation pointer only: the counted proof cites this statement by \textbackslash{}ref as a standing input; the deduction the counted proof has to make is carried out inside the counted span, and no part of that deduction is delegated to this pointer.\par\medskip
\ifdefined\hypertarget\hypertarget{sourcepointer-4}{}\fi
\noindent\textit{Source reference.} Original reference 04SI: focused-seeds/geometry/S9-alg-representable-diagonal/ref-tags.txt, source lines 6--6. Navigation pointer only: the counted proof cites this statement by \textbackslash{}ref as a standing input; the deduction the counted proof has to make is carried out inside the counted span, and no part of that deduction is delegated to this pointer.\par\medskip
\begin{lemma}
\label{lemma-representable-diagonal}
Let $S$ be a scheme contained in $\Sch_{fppf}$.
Let $\mathcal{X}$ be a category fibred in groupoids over
$(\Sch/S)_{fppf}$. The following are equivalent:
\begin{enumerate}
\item the diagonal $\mathcal{X} \to \mathcal{X} \times \mathcal{X}$
is representable by algebraic spaces,
\item for every scheme $U$ over $S$, and any
$x, y \in \Ob(\mathcal{X}_U)$ the sheaf
$\mathit{Isom}(x, y)$ is an algebraic space over $U$,
\item for every scheme $U$ over $S$, and any $x \in \Ob(\mathcal{X}_U)$
the associated $1$-morphism $x : (\Sch/U)_{fppf} \to \mathcal{X}$
is representable by algebraic spaces,
\item for every pair of schemes $T_1, T_2$ over $S$, and any
$x_i \in \Ob(\mathcal{X}_{T_i})$, $i = 1, 2$ the $2$-fibre product
$(\Sch/T_1)_{fppf} \times_{x_1, \mathcal{X}, x_2}
(\Sch/T_2)_{fppf}$
is representable by an algebraic space,
\item for every representable category fibred in groupoids $\mathcal{U}$
over $(\Sch/S)_{fppf}$ every $1$-morphism
$\mathcal{U} \to \mathcal{X}$ is representable by algebraic spaces,
\item for every pair $\mathcal{T}_1, \mathcal{T}_2$ of representable
categories fibred in groupoids over $(\Sch/S)_{fppf}$ and any
$1$-morphisms $x_i : \mathcal{T}_i \to \mathcal{X}$, $i = 1, 2$ the
$2$-fibre product $\mathcal{T}_1 \times_{x_1, \mathcal{X}, x_2} \mathcal{T}_2$
is representable by an algebraic space,
\item for every category fibred in groupoids $\mathcal{U}$
over $(\Sch/S)_{fppf}$ which is
representable by an algebraic space every $1$-morphism
$\mathcal{U} \to \mathcal{X}$ is representable by algebraic spaces,
\item for every pair $\mathcal{T}_1, \mathcal{T}_2$ of categories fibred
in groupoids over $(\Sch/S)_{fppf}$ which are representable
by algebraic spaces, and any $1$-morphisms
$x_i : \mathcal{T}_i \to \mathcal{X}$ the
$2$-fibre product $\mathcal{T}_1 \times_{x_1, \mathcal{X}, x_2} \mathcal{T}_2$
is representable by an algebraic space.
\end{enumerate}
\end{lemma}

\begin{proof}
The equivalence of (1) and (2) follows from
Stacks, Lemma \ref{stacks-lemma-isom-as-2-fibre-product}
and the definitions.
Let us prove the equivalence of (1) and (3).
Write $\mathcal{C} = (\Sch/S)_{fppf}$ for the base category.
We will use some of the observations of the proof of the similar
Categories, Lemma \ref{categories-lemma-representable-diagonal-groupoids}.
We will use the symbol $\cong$ to mean ``equivalence of categories fibred
in groupoids over $\mathcal{C} = (\Sch/S)_{fppf}$''.
Assume (1). Suppose given $U$ and $x$ as in (3). For any scheme $V$
and $y \in \Ob(\mathcal{X}_V)$ we see (compare reference above) that
$$
\mathcal{C}/U
\times_{x, \mathcal{X}, y}
\mathcal{C}/V
\cong
(\mathcal{C}/U \times_S V)
\times_{(x, y), \mathcal{X} \times \mathcal{X}, \Delta}
\mathcal{X}
$$
which is representable by an algebraic space by assumption. Conversely,
assume (3). Consider any scheme $U$ over $S$ and a pair $(x, x')$
of objects of $\mathcal{X}$ over $U$. We have to show that
$\mathcal{X} \times_{\Delta, \mathcal{X} \times \mathcal{X}, (x, x')} U$
is representable by an algebraic space. This is clear because
(compare reference above)
$$
\mathcal{X}
\times_{\Delta, \mathcal{X} \times \mathcal{X}, (x, x')}
\mathcal{C}/U
\cong
(\mathcal{C}/U \times_{x, \mathcal{X}, x'} \mathcal{C}/U)
\times_{\mathcal{C}/U \times_S U, \Delta}
\mathcal{C}/U
$$
and the right hand side is representable by an algebraic space by assumption
and the fact that the category of algebraic spaces over $S$ has fibre products
and contains $U$ and $S$.

\medskip\noindent
The equivalences
(3) $\Leftrightarrow$ (4),
(5) $\Leftrightarrow$ (6),
and
(7) $\Leftrightarrow$ (8)
are formal. The equivalences
(3) $\Leftrightarrow$ (5) and
(4) $\Leftrightarrow$ (6)
follow from
Lemma \ref{lemma-representable-by-spaces-morphism-equivalent}.
Assume (3), and let $\mathcal{U} \to \mathcal{X}$ be as in (7).
To prove (7) we have to show that for every scheme $V$ and $1$-morphism
$y : (\Sch/V)_{fppf} \to \mathcal{X}$ the $2$-fibre product
$(\Sch/V)_{fppf} \times_{y, \mathcal{X}} \mathcal{U}$
is representable by an algebraic space. Property (3) tells us
that $y$ is representable by algebraic spaces hence
Lemma \ref{lemma-base-change-by-space-representable-by-space}
implies what we want. Finally, (7) directly implies (3).
\end{proof}

\end{document}
